\chapter{Análise de Vivacidade e Alocação de
Registradores}\label{chap:x86-alloc}

\section{Introdução}

A IRT canônica produzida pelo capítulo anterior contém temporários em
quantidade potencialmente ilimitada. O hardware X86-64 oferece apenas
dez registradores no pool de alocação. A \textbf{alocação de registradores}
é o processo de decidir, para cada temporário, se ele vai residir num
registrador físico ou num slot na pilha (o chamado \emph{derramamento} ou
\emph{spilling}).

A chave para uma alocação correta é saber quais temporários precisam
existir simultaneamente. Dois temporários que nunca estão ``vivos'' ao
mesmo tempo podem compartilhar o mesmo registrador sem problema. A
\textbf{análise de vivacidade} (\textit{liveness analysis}) computa
exatamente esse conjunto para cada ponto do programa.

Com a informação de vivacidade em mãos, o problema de alocação se
reduz a $k$-coloração de grafos: os temporários são nós, há uma aresta
entre dois temporários que interferem (estão vivos ao mesmo tempo) e
devem receber a mesma ``cor'' (registrador). O clássico algoritmo de
Kempe oferece uma solução eficiente, baseada em simplificação iterativa
do grafo.

Este capítulo cobre:

\begin{enumerate}
  \item Definição formal de vivacidade e as equações de ponto fixo que
    a calculam.
  \item Construção do grafo de interferência a partir da vivacidade.
  \item O algoritmo de Kempe: simplificação, escolha de candidato a
    derramamento e seleção de cores.
  \item Implementação completa nos módulos \texttt{Liveness},
    \texttt{IGraph} e \texttt{GraphColor}.
\end{enumerate}

\section{Análise de Vivacidade}

\subsection{Definição}

\begin{Definition}[Vivacidade]
  Um temporário $t$ é \textbf{vivo} num ponto $p$ do programa se existe
  algum caminho de $p$ até um \emph{uso} de $t$ que não passa por
  nenhuma \emph{definição} de $t$.
\end{Definition}

Intuitivamente, $t$ é vivo em $p$ quando seu valor atual ``ainda vai ser
necessário'' em alguma execução futura possível. Se $t$ não está vivo em
$p$, seu registrador pode ser reutilizado por outro temporário.

\subsection{Equações de Ponto Fixo}

Dado o CFG linearizado com instruções $s_0, s_1, \ldots, s_{n-1}$,
definimos para cada índice $i$:

\begin{align*}
  \mathit{use}[i] &= \text{conjunto de temporários \emph{lidos} por } s_i \\
  \mathit{def}[i] &= \text{conjunto de temporários \emph{escritos} por } s_i
\end{align*}

As equações de \textbf{análise retroativa} são:

\[
  \mathit{live\_out}[i] = \bigcup_{j \in \mathit{succs}(i)} \mathit{live\_in}[j]
\]

\[
  \mathit{live\_in}[i] = \mathit{use}[i] \cup
  \bigl(\mathit{live\_out}[i] \setminus \mathit{def}[i]\bigr)
\]

A análise itera essas equações partindo de $\mathit{live\_in}[i] = \emptyset$
para todo $i$, até que nenhum conjunto se altere (ponto fixo). A solução
é a \emph{menor} solução: o conjunto mais restrito de temporários que
satisfaz as equações.

\begin{remark}
  A análise é \emph{retroativa} porque vivacidade é uma propriedade do
  futuro de um ponto de execução: para saber se $t$ é vivo antes de
  $s_i$, precisamos saber se $t$ é usado depois de $s_i$: informação
  que está nos \emph{sucessores} de $i$ no CFG. Por isso, iterar de
  trás para frente dentro de cada passagem (ordem de Gauss-Seidel)
  acelera a convergência.
\end{remark}

\subsection{Use e Def para Instruções IRT}

A tabela a seguir define $\mathit{use}[i]$ e $\mathit{def}[i]$ para
cada forma de enunciado IRT após a canonicalização:

\begin{center}
  \begin{tabular}{lll}
    \toprule
    Enunciado $s_i$ & $\mathit{use}[i]$ & $\mathit{def}[i]$ \\
    \midrule
    \texttt{MOVE(TEMP t, e)}      & $\mathit{temps}(e)$
    & $\{t\}$   \\
    \texttt{MOVE(MEM addr, val)}  & $\mathit{temps}(\mathit{addr})
    \cup \mathit{temps}(\mathit{val})$ & $\emptyset$ \\
    \texttt{EXP e}                & $\mathit{temps}(e)$
    & $\emptyset$ \\
    \texttt{CJUMP e \_ \_}       & $\mathit{temps}(e)$
    & $\emptyset$ \\
    \texttt{JUMP e}               & $\mathit{temps}(e)$
    & $\emptyset$ \\
    \texttt{RETURN [e]}           & $\mathit{temps}(e)$
    & $\emptyset$ \\
    \texttt{LABEL \_}             & $\emptyset$
    & $\emptyset$ \\
    \bottomrule
  \end{tabular}
\end{center}

Onde $\mathit{temps}(e)$ coleta recursivamente todos os \texttt{TEMP}
em $e$. Apenas \texttt{MOVE(TEMP t, \_)} constitui uma \emph{definição}
de $t$; todas as demais ocorrências de temporários, incluindo
\texttt{MOVE(MEM addr, val)} onde $addr$ e $val$ são lidos mas nenhum
temporário é escrito, são usos.

\subsection{Implementação}

O módulo \texttt{IR.Backend.X86.Liveness} implementa as funções de
uso e definição:

\begin{minted}{haskell}
freeTempsExpr :: Expr -> Set Temp
freeTempsExpr (TEMP t)        = Set.singleton t
freeTempsExpr (BINOP _ e1 e2) = freeTempsExpr e1 `Set.union` freeTempsExpr e2
freeTempsExpr (MEM e)         = freeTempsExpr e
freeTempsExpr (CALL ef args)  = Set.unions (map freeTempsExpr (ef : args))
freeTempsExpr (ESEQ s e)      = useTemp s `Set.union` freeTempsExpr e
freeTempsExpr _               = Set.empty  -- CONST, NAME

useTemp :: Stmt -> Set Temp
useTemp (MOVE (TEMP _)   e)   = freeTempsExpr e
useTemp (MOVE dst        src) = freeTempsExpr dst `Set.union` freeTempsExpr src
useTemp (EXP e)               = freeTempsExpr e
useTemp (RETURN es)           = Set.unions (map freeTempsExpr es)
useTemp (CJUMP e _ _)         = freeTempsExpr e
useTemp _                     = Set.empty

defTemp :: Stmt -> Set Temp
defTemp (MOVE (TEMP t) _) = Set.singleton t
defTemp _                 = Set.empty
\end{minted}

E o cálculo de ponto fixo:

\begin{minted}{haskell}
type LiveMap = Map Int (Set Temp)   -- live_in[i]

liveness :: FuncCFG -> LiveMap
liveness cfg = fixpoint initMap
  where
    idxs    = reverse (cfgIndices cfg)  -- ordem retroativa
    initMap = Map.fromList [(i, Set.empty) | i <- cfgIndices cfg]

    fixpoint m =
      let m' = foldl step m idxs
      in if m' == m then m else fixpoint m'

    step acc i =
      let lo    = Set.unions [acc Map.! j | j <- succs cfg i]
          newIn = useTemp (cfgStmtAt cfg i)
                  `Set.union`
                  (lo `Set.difference` defTemp (cfgStmtAt cfg i))
      in Map.insert i newIn acc
\end{minted}

A função \texttt{liveOut cfg liveIn i} calcula $\mathit{live\_out}[i]$
a partir do \texttt{LiveMap} finalizado:

\begin{minted}{haskell}
liveOut :: FuncCFG -> LiveMap -> Int -> Set Temp
liveOut cfg liveIn i =
  Set.unions [liveIn Map.! j | j <- succs cfg i]
\end{minted}

\section{Grafo de Interferência}

\subsection{Definição}

\begin{Definition}[Interferência]
  Dois temporários $t$ e $u$ \textbf{interferem} se existe alguma
  instrução $i$ tal que $t \in \mathit{def}[i]$ e
  $u \in \mathit{live\_out}[i]$ (ou vice-versa), com $t \neq u$.
\end{Definition}

Temporários que interferem não podem ocupar o mesmo registrador, pois na
instrução de definição ambos precisam de locais distintos: $t$ está sendo
criado e $u$ ainda será necessário.

O \textbf{grafo de interferência} $G = (V, E)$ tem:

\begin{itemize}
  \item $V$ = conjunto de todos os temporários do programa.
  \item $(t, u) \in E$ se e somente se $t$ e $u$ interferem.
\end{itemize}

A regra de construção é:

\[
  \infer{\text{adicionar aresta } (t, u) \text{ a } G}
  {t \in \mathit{def}[i] \quad u \in \mathit{live\_out}[i] \quad t \neq u}
\]

\subsection{Parâmetros com Interferência Mútua}

Os parâmetros formais de uma função chegam carregados em registradores
distintos da ABI (\texttt{\%rdi}, \texttt{\%rsi}, etc.) e estão todos
vivos \emph{simultaneamente} no ponto de entrada da função. A análise de
vivacidade pode não ver isso explicitamente — se um parâmetro não possui
uso após a instrução de entrada, ele pode não aparecer em nenhum
$\mathit{live\_in}[i]$.

Por isso, antes de iniciar a coloração, adicionamos arestas entre todos
os pares de parâmetros formais:

\[
  \forall\, p_1, p_2 \in \mathit{params},\; p_1 \neq p_2 \implies
  (p_1, p_2) \in E
\]

Sem essa inserção, dois parâmetros distintos podem ser mapeados ao mesmo
registrador — um erro silencioso que produz resultados errados.

\subsection{Estrutura de Dados}

O módulo \texttt{IR.Backend.X86.IGraph} implementa o grafo de
interferência com lista de adjacência e cache de graus:

\begin{minted}{haskell}
data IGraph = IGraph
  { igAdj :: Map Temp (Set Temp)   -- adjacências (simétrica)
  , igDeg :: Map Temp Int           -- cache de graus
  } deriving (Eq, Show)

addEdge :: Temp -> Temp -> IGraph -> IGraph
addEdge t u g
  | t == u    = addNode t g
  | u `Set.member` neighbors t g' = g'
  | otherwise = g'
      { igAdj = Map.adjust (Set.insert u) t
              $ Map.adjust (Set.insert t) u
              $ igAdj g'
      , igDeg = Map.adjust (+1) t
              $ Map.adjust (+1) u
              $ igDeg g'
      }
  where g' = addNode u (addNode t g)

removeNode :: Temp -> IGraph -> IGraph
removeNode t g =
  let nbrs = neighbors t g
      g'   = g { igAdj = Map.delete t (igAdj g)
               , igDeg = Map.delete t (igDeg g) }
  in Set.foldl' (\acc u ->
       acc { igAdj = Map.adjust (Set.delete t) u (igAdj acc)
           , igDeg = Map.adjust (subtract 1)   u (igDeg acc) })
     g' nbrs
\end{minted}

O cache \texttt{igDeg} evita recomputar o grau de cada nó a cada passo
da simplificação, o que seria $O(n)$ por nó sem o cache.

\subsection{Construção do Grafo}

\begin{minted}{haskell}
buildIGraph :: FuncCFG -> LiveMap -> IGraph
buildIGraph cfg liveIn =
  let allTemps =
        Set.unions (Map.elems liveIn)
        `Set.union`
        Set.unions [defTemp (cfgStmtAt cfg i) | i <- cfgIndices cfg]
      g0 = Set.foldl' (flip addNode) emptyIGraph allTemps
  in foldl' addInstr g0 (cfgIndices cfg)
  where
    addInstr g i =
      let tSet = defTemp (cfgStmtAt cfg i)
          out  = liveOut cfg liveIn i
      in if Set.null tSet then g
         else Set.foldl' (\g' u -> addEdge (Set.findMin tSet) u g') g out
\end{minted}

Todos os temporários são adicionados como nós, mesmo aqueles sem arestas
(temporários cujo valor é definido mas nunca está vivo ao mesmo tempo que
outro temporário). Isso garante que recebam uma cor mesmo sem restrições.

\section{Alocação por Coloração de Grafo}

\subsection{O Problema}

Dado um grafo de interferência $G$ e um conjunto de $k$ ``cores''
(registradores do pool de alocação), o problema é atribuir uma cor a
cada temporário de modo que temporários adjacentes recebam cores
distintas. Se não for possível colorir algum nó com as cores
disponíveis, ele é \textbf{derramado} (\textit{spilled}): seu valor é
armazenado na pilha e carregado/guardado a cada uso/definição.

O pool de alocação tem $k = 10$ registradores:
\texttt{\%rbx}, \texttt{\%rcx}, \texttt{\%rsi}, \texttt{\%rdi},
\texttt{\%r8}, \texttt{\%r9}, \texttt{\%r12}–\texttt{\%r15}.

\subsection{Preferência de Cores: Callee-Saved Primeiro}

As cores são oferecidas ao nó em uma ordem de preferência:
\emph{callee-saved primeiro}, depois os caller-saved do pool.

\[
  \mathit{preferOrder} = [\texttt{\%rbx},\, \texttt{\%r12},\, \texttt{\%r13},\,
    \texttt{\%r14},\, \texttt{\%r15},\, \texttt{\%rcx},\, \texttt{\%rsi},\,
  \texttt{\%rdi},\, \texttt{\%r8},\, \texttt{\%r9}]
\]

A razão: um temporário que recebe um registrador callee-saved e sobrevive
a uma chamada de função não precisa ser salvo e restaurado em torno do
\texttt{callq} — o próprio protocolo de chamada garante isso. Já um
temporário em registrador caller-saved precisaria ser explicitamente salvo
antes de cada chamada e restaurado depois.

\subsection{Algoritmo de Kempe}

O algoritmo opera em duas fases principais.

\subsubsection{Fase 1 — Simplificação}

Enquanto houver nós com grau $< k$, remova um deles e empilhe-o.
Se todos os nós tiverem grau $\geq k$, escolha um candidato a
derramamento e remova-o (marcando-o como potencial derramamento).

A propriedade fundamental é:

\begin{Theorem}[Simplificação segura]
  Se $G$ tem um nó $t$ com $\mathit{grau}(t) < k$, então $G$ é
  $k$-colorível se e somente se $G \setminus \{t\}$ é $k$-colorível.
\end{Theorem}

\emph{Demonstração.} Se $G \setminus \{t\}$ é $k$-colorível, os
vizinhos de $t$ usam no máximo $k-1$ cores distintas, logo existe
sempre pelo menos uma cor livre para $t$. \(\square\)

Isso significa que podemos simplificar com segurança qualquer nó de
grau $< k$: a sua remoção não pode transformar um grafo colorível em
não-colorível.

Para a escolha do candidato a derramamento quando não há nó de grau
$< k$, usamos uma heurística: derrama o temporário com o
\textbf{maior intervalo de vivacidade}.

\begin{Definition}[Intervalo de vivacidade]
  O intervalo de $t$ é $[\mathit{start}(t), \mathit{end}(t)]$ onde
  $\mathit{start}(t)$ é o menor índice em que $t$ está vivo ou é
  definido, e $\mathit{end}(t)$ é o maior índice em que $t$ está vivo.
  O comprimento é $\mathit{end}(t) - \mathit{start}(t)$.
\end{Definition}

Um temporário com intervalo longo exerce pressão sobre o pool por
mais instruções, sendo o candidato mais natural a ceder seu registrador.

\subsubsection{Fase 2 — Seleção de Cores}

Após o empilhamento, desempilhamos os temporários um a um e atribuímos
a menor cor disponível (na ordem de preferência). ``Disponível'' significa
não usada por nenhum vizinho \emph{no grafo original} (antes das remoções).

\[
  \infer{\text{atribuir } c^* = \min(\mathit{Pool} \setminus
  C_{\mathit{usadas}}) \text{ a } t}
  {C_{\mathit{usadas}} = \{c \mid (t, u) \in E_{\mathit{orig}},\;
  \mathit{cor}(u) = c\}}
\]

Se $\mathit{Pool} \setminus C_{\mathit{usadas}} = \emptyset$, o nó é
derramado: recebe um offset de memória negativo em relação a
\texttt{\%rbp} ($-8$, $-16$, etc.).

A fase de seleção usa o grafo \emph{original} (não o simplificado) para
consultar os vizinhos. As remoções da fase 1 são apenas um dispositivo
para determinar a \emph{ordem} de coloração, não para alterar as restrições.

\subsection{Implementação: Simplificação}

\begin{minted}{haskell}
simplify :: IGraph -> IntervalMap -> ([Temp], Set Temp)
simplify ig0 ivm = go ig0 [] Set.empty (Set.toList (nodes ig0))
  where
    go _  stk spills []  = (stk, spills)
    go ig stk spills wl  =
      case findLow ig wl of
        Just t  ->
          go (removeNode t ig) (t : stk) spills (filter (/= t) wl)
        Nothing ->
          let t = maximumBy
                    (comparing (\t' -> maybe 0 ivLen (Map.lookup t' ivm)))
                    wl
          in go (removeNode t ig) (t : stk) (Set.insert t spills)
                (filter (/= t) wl)

findLow :: IGraph -> [Temp] -> Maybe Temp
findLow ig = listToMaybe . filter (\t -> degree t ig < k)
\end{minted}

\subsection{Implementação: Seleção}

\begin{minted}{haskell}
selectColors :: [Temp] -> Set Temp -> IGraph -> Int -> (Coloring, Int)
selectColors stack _potSpills origGraph startSlot =
  foldl' step (Map.empty, startSlot) stack
  where
    step (col, nextSlot) t =
      let used  = Set.fromList
                    [ r | u <- Set.toList (neighbors t origGraph)
                        , Just (InReg r) <- [Map.lookup u col] ]
          avail = filter (`Set.notMember` used) preferOrder
      in case avail of
           (r : _) -> (Map.insert t (InReg r)       col, nextSlot)
           []      -> (Map.insert t (InMem nextSlot) col, nextSlot - 8)
\end{minted}

O resultado é um \texttt{Coloring} (mapa de temporário para
\texttt{TempLoc}) e o próximo offset de spill disponível. Slots de spill
são alocados em $-8, -16, -24, \ldots$ a partir de \texttt{\%rbp}.

\subsection{Ponto de Entrada: \texttt{colorFunc}}

\begin{minted}{haskell}
colorFunc :: [Temp] -> Stmt -> (Coloring, Int)
colorFunc params body =
  let cfg        = buildCFG body
      livm       = liveness cfg
      ivm        = buildIntervalMap cfg livm
      ig0        = buildIGraph cfg livm
      ig         = addParamEdges params ig0
      (stack, potSpills) = simplify ig ivm
      (col, nextSlot)    = selectColors stack potSpills ig (-8)
      spillBytes         = negate nextSlot - 8
  in (col, spillBytes)
  where
    addParamEdges ps g =
      foldr (\(p1, p2) acc -> addEdge p1 p2 acc) g
            [(p1, p2) | p1 <- ps, p2 <- ps, p1 /= p2]
\end{minted}

O valor \texttt{spillBytes} é o número de bytes a reservar na pilha para
temporários derramados ($0$ quando não há derramamento).

\section{Layout de Frame e Alinhamento}

\subsection{Layout}

O frame de ativação de uma função tem a seguinte estrutura (crescendo
de cima para baixo, em direção a endereços menores):

\begin{center}
  \begin{tabular}{|c|}
    \hline
    \ldots (frame do chamador) \\
    \hline
    endereço de retorno \quad [\texttt{callq} empilha] \\
    \hline
    \texttt{\%rbp} antigo \quad [\texttt{pushq \%rbp}] \\
    \hline
    registradores callee-saved \quad [$N \times 8$ bytes] \\
    \hline
    slots de spill \quad [$\mathit{spillBytes}$ bytes] \\
    \hline
    padding de alinhamento \\
    \hline
  \end{tabular}
\end{center}

\subsection{Cálculo de \texttt{frameAdj}}

O ABI System V AMD64 exige que \texttt{\%rsp} seja múltiplo de 16 bytes
no momento da instrução \texttt{callq} (que empilha 8 bytes). Após o
\texttt{pushq \%rbp} o stack está alinhado a 16. Cada push de registrador
callee-saved desloca 8 bytes. Precisamos de um ajuste $F \geq
\mathit{spillBytes}$
tal que:

\[
  8N + F \equiv 0 \pmod{16}
\]

onde $N$ é o número de registradores callee-saved empilhados no prólogo.

\begin{minted}{haskell}
frameAdj :: Int -> Int -> Int
frameAdj nCallee spillBytes =
  let base     = align16 spillBytes
      total    = 8 * nCallee + base
      misalign = total `mod` 16
  in if misalign == 0 then base else base + (16 - misalign)

align16 :: Int -> Int
align16 n = ((n + 15) `div` 16) * 16
\end{minted}

Se $F = 0$, a instrução \texttt{subq \$0, \%rsp} é omitida.

\section{Exemplo: Alocação em \texttt{factorial}}

\subsection{IRT Canônica (simplificada)}

\begin{center}
  \begin{tabular}{ll}
    \toprule
    Índice & Enunciado \\
    \midrule
    0 & \texttt{MOVE(TEMP n, TEMP arg)} \\
    1 & \texttt{MOVE(TEMP t0, BINOP(EQ, TEMP n, CONST 0))} \\
    2 & \texttt{CJUMP(TEMP t0, .true, .false)} \\
    3 & \texttt{LABEL .true} \\
    4 & \texttt{RETURN [CONST 1]} \\
    5 & \texttt{LABEL .false} \\
    6 & \texttt{MOVE(TEMP t1, BINOP(SUB, TEMP n, CONST 1))} \\
    7 & \texttt{MOVE(TEMP r, CALL(factorial, [TEMP t1]))} \\
    8 & \texttt{MOVE(TEMP t2, BINOP(MUL, TEMP n, TEMP r))} \\
    9 & \texttt{RETURN [TEMP t2]} \\
    \bottomrule
  \end{tabular}
\end{center}

\subsection{Vivacidade}

\begin{center}
  \begin{tabular}{cllll}
    \toprule
    $i$ & $\mathit{use}[i]$ & $\mathit{def}[i]$ &
    $\mathit{live\_in}[i]$ & $\mathit{live\_out}[i]$ \\
    \midrule
    0 & \{arg\}    & \{n\}  & \{arg\}    & \{n\}       \\
    1 & \{n\}      & \{t0\} & \{n\}      & \{n, t0\}   \\
    2 & \{t0\}     & $\emptyset$ & \{n, t0\} & \{n\}  \\
    4 & $\emptyset$ & $\emptyset$ & $\emptyset$ & $\emptyset$ \\
    6 & \{n\}      & \{t1\} & \{n\}      & \{n, t1\}   \\
    7 & \{t1\}     & \{r\}  & \{n, t1\}  & \{n, r\}    \\
    8 & \{n, r\}   & \{t2\} & \{n, r\}   & \{t2\}      \\
    9 & \{t2\}     & $\emptyset$ & \{t2\} & $\emptyset$ \\
    \bottomrule
  \end{tabular}
\end{center}

\subsection{Grafo de Interferência}

As arestas são adicionadas pela regra: $t \in \mathit{def}[i]$,
$u \in \mathit{live\_out}[i]$, $t \neq u$:

\begin{itemize}
  \item $i=1$: def=\{t0\}, live\_out=\{n, t0\} $\to$ aresta $(t0, n)$
  \item $i=6$: def=\{t1\}, live\_out=\{n, t1\} $\to$ aresta $(t1, n)$
  \item $i=7$: def=\{r\}, live\_out=\{n, r\} $\to$ aresta $(r, n)$
  \item $i=8$: def=\{t2\}, live\_out=\{t2\} $\to$ nenhuma (t2 não
    interfere com outros)
\end{itemize}

Arestas de parâmetros: como \texttt{factorial} tem um único parâmetro
\texttt{arg}, não há par para inserir aresta mútua.

\medskip
\noindent
Resumo do grafo:
$n \leftrightarrow t0$, $n \leftrightarrow t1$, $n \leftrightarrow r$.
O nó $n$ tem grau 3; os demais têm grau $\leq 1$.

\subsection{Coloração}

$k = 10 > 3$, logo todos os nós têm grau $< k$. Ordem de simplificação
(grau do menor ao maior): \texttt{arg}, \texttt{t0}, \texttt{t1},
\texttt{r}, \texttt{t2}, \texttt{n}.

Seleção (desempilhando, preferência callee-saved):

\begin{center}
  \begin{tabular}{lll}
    \toprule
    Temporário & Vizinhos já coloridos & Cor atribuída \\
    \midrule
    \texttt{n}   & $\emptyset$                    & \texttt{\%rbx} \\
    \texttt{t2}  & $\emptyset$                    & \texttt{\%rbx} \\
    \texttt{r}   & \{n=\texttt{\%rbx}\}           & \texttt{\%r12} \\
    \texttt{t1}  & \{n=\texttt{\%rbx}\}           & \texttt{\%r12} \\
    \texttt{t0}  & \{n=\texttt{\%rbx}\}           & \texttt{\%r12} \\
    \texttt{arg} & $\emptyset$                    & \texttt{\%rbx} \\
    \bottomrule
  \end{tabular}
\end{center}

\texttt{n} e \texttt{t2} recebem \texttt{\%rbx} sem conflito porque seus
intervalos não se sobrepõem. Da mesma forma, \texttt{r}, \texttt{t1} e
\texttt{t0} recebem \texttt{\%r12} porque nunca estão todos vivos ao mesmo
tempo.

O prólogo empilha \texttt{\%rbx} e \texttt{\%r12} ($N = 2$).
$\mathit{spillBytes} = 0$.
$\mathit{frameAdj}(2, 0) = \mathit{align16}(0) = 0$, com
$8 \cdot 2 + 0 = 16 \equiv 0 \pmod{16}$: sem \texttt{subq}.

\section{Conclusão}

A análise de vivacidade e a alocação por coloração de grafos são as duas
fases que determinam a \emph{qualidade} do código gerado: um coloridor
eficiente minimiza os derramamentos, reduzindo os acessos à memória; um
cálculo correto de vivacidade garante que dois temporários vivos ao mesmo
tempo nunca compartilhem um registrador.

A adição explícita de arestas entre parâmetros formais é um detalhe
crucial que não aparece na formulação tradicional do algoritmo de Kempe,
mas que é indispensável quando os parâmetros chegam em registradores da
ABI e são copiados para temporários antes de serem usados.

No próximo capítulo traduzimos a IRT anotada com o \texttt{Coloring} para
instruções X86-64 concretas, gerenciando o alinhamento de pilha em
chamadas aninhadas e emitindo o prólogo e o epílogo corretos para cada função.
